\chapter{Abstract}


Despite girls achieving academic results that are comparable to their male peers in early secondary education, significant disparities still exist regarding their self-efficacy in mathematics, anxiety levels, and long-term STEM career choices. This dissertation addresses these systemic inequities by exploring the potential of narrative-driven educational practices aimed at challenging the traditional, male-dominated perception of mathematical genius.

The dissertation combines a historical reconstruction of women's contributions to mathematics—from antiquity to the twentieth century, with particular attention to Maryam Mirzakhani—with an empirical field study assessing the short-term effects of a narrative-based school intervention.

In December 2025, a school-based quasi-experimental intervention took place in Trieste, involving a sample of 362 lower secondary school students across three grade levels. Pre- and post-test standardized questionnaires were analyzed using Jamovi to evaluate short-term changes in math anxiety, mathematics self-efficacy, attitudes toward the subject, and explicit endorsement of gender stereotypes.

The intervention was associated with a marked reduction in explicit gender stereotypes and a significant decrease in mathematics anxiety, alongside a slight improvement in students' general attitudes toward mathematics. By contrast, mathematics self-efficacy remained stable, suggesting that changes in competence beliefs require more sustained mastery experiences.

Ultimately, this study demonstrates that integrating historical female role models into mathematics education is a promising strategy for challenging exclusionary representations of the discipline and fostering a more inclusive understanding of mathematics as an accessible, creative, and universally human endeavor.
